\subsection{Genus-zero generation of abelian covers and elliptic Galois covers}

\begin{claim}{FF.48}{Proved}\label{clm:FFlatteskernel}
Let \(E/B_p\) be finite, and let \(\varphi\) be a pinned
automorphism of \((\mathbb A^1_E)_{\acute et}\) over \(BG_E\)
whose actual geometric point function is the identity.
The natural scheme realizations of
Claim~\ref{clm:FFgenuszerokernel} extend to the following
objects, after every finite extension of constants:
finite fibre products of the genus-zero objects, their
open subobjects, and finite etale quotients by finite
groups. In particular they extend to every subcover of
the Galois closure of a connected finite etale
genus-zero cover of an open of \(\mathbb A^1\).

They consequently extend to every etale restriction of
a finite abelian cover of \(\mathbb P^1_E\), of
arbitrary genus, and to every finite Galois cover
\[
                  \pi:C\longrightarrow\mathbb P^1_E
       \quad\hbox{with }g(C)=1.
                                                        \tag{FF.48.1}
\]
Here ``extend'' means that the induced slice point
function is realized by a unique scheme isomorphism
\(\varphi^*V\to V\) over the chosen base open, and
these isomorphisms are natural in every morphism between
the stated objects. The statement includes a quadratic
cover with four geometric branch points. It also
includes every degree-three cover with genus-one
source, without a Galois assumption.

Here \(C\) is smooth and geometrically integral.
Thus, after identifying the inverse image of any
connected cyclic torsor of the stated kind with the
original one, its possible fibrewise deck rotation is
a single global deck transformation. Point-dependent
rotations on cyclic covers of arbitrary genus are
excluded. No generation or cofinality
claim for all positive-genus covers is asserted.
\end{claim}

\begin{proof}
All constructions below may be made after removing
finitely many points of the base. The resulting objects
are quasicompact separated etale objects, and their
inverse images are represented by objects of the same
kind. We keep the structural comparisons over the
constant field throughout.

\emph{Closure under the stated constructions.}
Say that an etale object \(V\to X\), for an open
\(X\subseteq\mathbb A^1_E\), is realized if its
actual slice point function is induced by an
\(X\)-isomorphism
\[
                       a_V:\varphi^*V\xrightarrow{\sim}V.
                                                        \tag{FF.48.2}
\]
Such an isomorphism is unique: two maps between these
reduced finite-type schemes, with separated target, that
agree on all geometric points agree as scheme maps.
Moreover, for any arrow \(V\to W\) between realized
objects, the two maps \(\varphi^*V\to W\) formed
using (FF.48.2) agree on geometric points by the
naturality of point evaluation. They therefore agree
as scheme maps. Realizations are consequently natural
on every full subcategory on which they exist.

If \(V_1,\ldots,V_r\) are realized, their fibre
product is realized by the product of the maps
\(a_{V_i}\). Indeed inverse image preserves finite
limits, and a geometric point of the product is
determined by its projections. If \(W\subseteq V\)
is an open subobject of a realized object, the image of
\(\varphi^*W\) under \(a_V\) is precisely \(W\):
the two opens have the same geometric points. Restriction
of \(a_V\) gives \(a_W\). This includes open-and-closed
connected components. These conclusions use the actual
point-function comparison, rather than just invariance
of the isomorphism class of \(V\).

Suppose that a finite group \(H\) acts on a realized
finite etale object \(Z\to X\). Naturality for every
action arrow identifies the transported action with the
given action on \(Z\). Inverse image preserves its
coequalizer. The finite etale quotient therefore has
an induced isomorphism
\[
             \varphi^*(Z/H)\ \xrightarrow{\sim}\ Z/H.
                                                        \tag{FF.48.3}
\]
It realizes the slice point function, because the finite
etale surjection \(Z\to Z/H\) is surjective on
geometric points and that point function is natural.

For clarity, the Galois-closure application can be made
entirely inside these constructions. If \(V\to X\)
is connected finite etale of degree \(d\), take in
\(V^d_X\) the open-and-closed subscheme of ordered
\(d\)-tuples of pairwise distinct geometric sheets.
It is the frame torsor of \(V\). A connected component
\(Z\) is its Galois closure, with group the monodromy
image in \(S_d\). Once \(V\) is realized, its frame
torsor and each specified component are realized by
the preceding paragraph. Each finite subcover of
\(Z\) is a quotient \(Z/H\), so (FF.48.3) applies.
This proves the first assertion using
Claim~\ref{clm:FFgenuszerokernel}.

\emph{All abelian covers of the rational base.}
After a finite extension of constants, every connected
cyclic cover of degree \(m\) is a Kummer cover.
Indeed all \(m\)-th roots of unity are available,
and Kummer theory applies to its cyclic function-field
extension. After a further finite extension that splits
the divisor and extracts a constant root, write it as
\[
       y^m=f(T)=c^m\prod_{i=1}^r(T-a_i)^{e_i},
           \qquad c\ne0,\quad e_i\in\mathbb Z.
                                                        \tag{FF.48.4}
\]
Remove the finitely many \(a_i\), the original affine
infinity, and any additional omitted base points.
Each cover \(u_i^m=T-a_i\) of the resulting open
has genus-zero completion and is realized by
Claim~\ref{clm:FFgenuszerokernel}. Their finite fibre
product is therefore realized, as are its connected
components and finite group quotients.

More explicitly, the classes of the distinct
\(T-a_i\) are independent in
\(\bar E(T)^\times/(\bar E(T)^\times)^m\), as is
seen from their respective valuations at \(a_i\).
Their compositum is thus a connected Galois cover
with group \(\mu_m^r\). The map
\[
       (u_1,\ldots,u_r)\longmapsto
                     y=c\prod_i u_i^{e_i}
                                                        \tag{FF.48.5}
\]
identifies (FF.48.4) with its quotient by the kernel
of the character
\((\zeta_i)\mapsto\prod_i\zeta_i^{e_i}\).
The character is surjective when the original cover
is geometrically connected of degree \(m\). If
constant splitting produces several components, apply
this argument to each component and descend their
union. This proves the cyclic assertion. A finite
abelian Galois extension is a compositum of cyclic
subextensions, by the elementary structure theorem
for its finite abelian Galois group. Its cover is a
connected component of the corresponding finite fibre
product, so every abelian cover is realized as well.

\emph{The elliptic Latt\`es construction.}
We first give the quadratic case explicitly. After a
finite Galois extension \(L/E\), a branch point of
\(\pi\) has an \(L\)-rational point above it, which
we choose as the origin of the elliptic curve \(A=C_L\).
The involution of the double cover is then \([-1]\).
Enlarge \(L\) so that all points of \(A[3]\) are
rational. Write
\[
                 \pi:A\longrightarrow A/\{\pm1\}
                                      =\mathbb P^1_L.
\]
Multiplication by three commutes with \([-1]\), and
therefore induces a rational morphism of degree nine
\[
          f_3:\mathbb P^1_L\longrightarrow\mathbb P^1_L,
                 \qquad f_3\circ\pi=\pi\circ[3].
                                                        \tag{FF.48.6}
\]
Let \(B\) be the branch-value set of \(\pi\), and
put \(X=\mathbb P^1_L\setminus(B\cup\{\infty\})\),
with \(\infty\) interpreted in the originally chosen
affine base coordinate. The maps
\[
 \begin{split}
  Z&=(\pi\circ[3])^{-1}(X)\longrightarrow X,\\
  V&=f_3^{-1}(X)\longrightarrow X
 \end{split}
                                                        \tag{FF.48.7}
\]
are finite etale: \([3]\) is etale, and the branch
values of the composite lie in \(B\). The second
cover has genus-zero completion.

The first cover in (FF.48.7) is Galois, with its
eighteen deck transformations
\[
            P\longmapsto \epsilon P+t,
             \qquad t\in A[3],\quad\epsilon\in\{1,-1\}.
                                                        \tag{FF.48.8}
\]
They are distinct and their number equals the degree
\(2\cdot9\) of the map. Its group is
\(G=A[3]\rtimes\{\pm1\}\). The intermediate
cover \(V\) is the quotient by
\(H=\{\pm1\}\). The core of \(H\) in \(G\)
is trivial: conjugation of \([-1]\) by translation
by \(t\ne0\) gives \(P\mapsto-P+2t\), which
does not belong to \(H\). Thus \(Z\) is the
Galois closure of the genus-zero cover \(V\to X\).
The first part of the proof realizes \(Z\).
Its quotient by the normal translation subgroup
\(A[3]\) is the original cover
\[
               \pi^{-1}(X)\longrightarrow X,
                                                        \tag{FF.48.9}
\]
with quotient map supplied by \([3]\). It is therefore
realized as well. The finite group and its faithful
degree-nine permutation action are also checked in
\path{agent_F/F_metabelian_checks.py}.

\emph{The other cyclic genus-one covers.}
The same argument does not require the deck group to
have order two. Over a finite extension, write a
generator of the cyclic deck group as
\(P\mapsto h(P)+t\) on an elliptic model of \(C\).
If \(h=1\), the group consists of translations and
its quotient has genus one, contrary to (FF.48.1).
Thus \(h-1\) is a nonzero endomorphism of the elliptic
curve, hence a surjective isogeny. The generator has
a fixed point; choose it as origin, after a further
finite extension. The deck group \(H\) now consists
of group automorphisms of \(A\), and
\(A/H=\mathbb P^1\).

Choose an integer \(n\ge3\) such that
\[
                 n^2>\deg(h-1)
                  \quad\text{for every }1\ne h\in H,
                                                        \tag{FF.48.10}
\]
and enlarge the finite constant field to contain
\(A[n]\). The map \([n]\) commutes with \(H\)
and induces a degree-\(n^2\) map \(f_n\) on
\(A/H=\mathbb P^1\). The composite
\(\pi\circ[n]\) is Galois with group
\(A[n]\rtimes H\). The core of \(H\) is trivial.
Indeed, an element \(h\) in that core must remain
in \(H\) after conjugation by every translation
\(t\in A[n]\), so \((h-1)t=0\) for all such
\(t\). If \(h\ne1\), this would put \(n^2\)
distinct points in the kernel of an isogeny of degree
strictly less than \(n^2\), contradicting (FF.48.10).
After removing the branch values and the affine
infinity, the preceding Galois-closure argument applies
unchanged. Quotienting by \(A[n]\) recovers
\(\pi\).

\emph{Every Galois genus-one cover.}
Let now \(\pi:A\to\mathbb P^1\) be a Galois
cover with finite deck group \(G\), and genus one
source. Work after a finite extension so that the
origins, deck transformations and torsion used below
are defined. The translations in \(G\) form a normal
subgroup \(T\). Its quotient
\[
                  \lambda:A\longrightarrow B=A/T
\]
is an isogeny of degree \(d=|T|\). The group
\(G/T\) is cyclic: the linear-part map embeds it in
the finite group of origin-preserving automorphisms
of an elliptic curve, which embeds by its action on
an invariant differential in the multiplicative
constant field. This action on differentials is
faithful in characteristic zero. The residual action
on \(B\) cannot consist only of translations,
since \(B/(G/T)=\mathbb P^1\). Its cyclic generator
therefore has a fixed point. Choose that point as
origin of \(B\), and a preimage as origin of \(A\),
so that \(\lambda\) remains a homomorphism. Put
\(H=G/T\), now acting on \(B\) by group
automorphisms, and let
\(\pi_B:B\to B/H=\mathbb P^1\).

Choose \(N\) divisible by \(d\), large enough
that \(N^2>\deg(h-1)\) for all \(1\ne h\in H\).
The Latt\`es argument shows that
\(\pi_B\circ[N]_B\), on the branch complement,
is the Galois closure of a genus-zero cover and is
therefore realized. The dual isogeny
\(\lambda^\vee:B\to A\) satisfies
\(\lambda\circ\lambda^\vee=[d]_B\). Thus
\[
 \pi\circ\lambda^\vee\circ[N/d]_B
                       =\pi_B\circ[N]_B.
                                                        \tag{FF.48.11}
\]
The original cover \(\pi\) is a subcover of this
realized finite Galois cover. It is realized by the
quotient part of the proof. This establishes the
full genus-one Galois assertion. The construction alone
does not treat every non-Galois genus-one map; the
following argument also settles all degree-three maps.

\emph{Degree-three genus-one covers, without a Galois assumption.}
The same generated category also contains every degree-three
cover with genus-one source. We give the finite monodromy
argument explicitly. It is the genus-one case of the
classical trigonal construction described in
\citet[Section~9.1, Remark~9.13]{FF_RecillasRodriguez2003};
the finite verification below makes this consequence
independent of a convention about the genus in a
moduli-space statement.

Work over \(\overline{\mathbb Q}\). A connected
degree-three cover has geometric monodromy \(C_3\)
or \(S_3\). The first case is already abelian.
In the second case let \(a\) be the number of branch
transpositions and \(b\) the number of branch
three-cycles. Riemann--Hurwitz and the product-one
relation give
\[
          a+2b=6,\qquad a>0\text{ even},\qquad
          (a,b)\in\{(6,0),(4,1),(2,2)\}.
                                                        \tag{FF.48.12}
\]
Let \(q:S_4\twoheadrightarrow S_3\) be the action
on the three partitions of four letters into two
unordered pairs. Its kernel is the Klein four group.
Every branch tuple in (FF.48.12) that generates
\(S_3\) and has product one has a lift through
\(q\) that generates \(S_4\), has product one,
and preserves each permutation index. Here the index
of a permutation on \(d\) letters is \(d\) minus
its number of cycles.

This last assertion is a finite lemma verified by
complete enumeration in the supplied program. The
five nonidentity elements of \(S_3\) give
\(5^4+5^5+5^6=19375\) tuples of the possible
lengths. Filtering by index sum six, product one,
and generation of \(S_3\) leaves respectively
72, 270, and 240 tuples. A transposition has two
index-preserving lifts through \(q\), and a
three-cycle has four. For every one of the 582
retained tuples the program enumerates all 64
possible index-preserving lifts and finds exactly
twelve with product one and generated group
\(S_4\). It verifies the quotient map by all
\(24^2\) group products and saves an explicit
lift for every retained tuple. Thus the finite
assertion checks all cases, not just selected
ramification patterns. The program and certificates
are supplied at the following paths.

\path{agent_F/F_metabelian_checks.py}

\path{agent_F/F_recillas_certificates.json}

Apply the characteristic-zero Riemann existence
theorem to a chosen lift at the actual branch
points of the original cover. Its degree-four
source has
\(2g-2=-2\cdot4+6=-2\), hence genus zero.
Its Galois closure has group \(S_4\). The quotient
through \(q\) has exactly the original
\(S_3\) branch tuple, so the original cubic
cover is the subcover associated with the inverse
image of a point stabilizer in \(S_3\).
It therefore belongs to the category already
realized above.

Here descent from the complex construction retains the
covering diagram, not merely its branch data. Let
\(X=\mathbb P^1_{\overline{\mathbb Q}}\setminus B\),
where \(B\) is the actual finite branch set, enlarged
by the original affine infinity if necessary. The
Riemann existence theorem identifies its complex finite
etale covers with finite topological covering spaces
\citep[Expose~XII, Theoreme~5.1]{FF_SGA1}.
Base change from \(\overline{\mathbb Q}\) to
\(\mathbb C\) is an equivalence on finite etale
covering categories: apply
\citet[Expose~XIII, Proposition~4.6(a)]{FF_SGA1}
with \(k=\overline{\mathbb Q}\), the smooth curve
\(X\), and \(Y=\operatorname{Spec}\mathbb C\).
The required resolution hypothesis holds for curves,
and in characteristic zero the prime-to-characteristic
quotient is the entire fundamental group. Since
\(\pi_1(Y)=1\), the resulting fundamental-group map
is an isomorphism. This equivalence is fully faithful,
so the \(S_4\)-action, quotient diagram and isomorphism
with the original cubic cover all descend to
\(\overline{\mathbb Q}\). They use finitely many
algebraic coefficients and morphisms, and therefore
descend to one finite extension \(E'/E\).
Uniqueness of the scheme realization with the specified
actual point function then descends that realization
to \(E\), as detailed below.

\emph{Descent and the fibre rotations.}
All the constructions just used occur after one finite
Galois extension \(L/E\). The resulting realization
of the original etale object is unique from its
geometric point function. That function is
\(G_E\)-equivariant, so the realization is fixed by
\(\operatorname{Gal}(L/E)\) and descends to \(E\).
If the originally specified etale restriction omits
further points of \(C\), use the open-subobject part
of the proof. Naturality follows from the uniqueness
argument at the beginning.

In any connected cyclic torsor case, compare the descended
realization with any chosen isomorphism of the
transported torsor with the original one. Their
composite is an automorphism of the connected cyclic
cover over the base, hence a single global deck
transformation. Accordingly the putative roots of unity
\(c_x\) are constant. This conclusion is about
the slice point function on this object; it does not
identify the induced equivalence on all its further
positive-genus covers.
\end{proof}


The generated subcategory in Claim~\ref{clm:FFlatteskernel}
is in fact not cofinal among all geometric covers.
By the proved fixed-genus composition-factor theorem,
there is a finite set \(\mathcal S_0\) of
nonalternating nonabelian simple groups such that every
simple composition factor of a genus-zero monodromy
group is cyclic, alternating, or belongs to
\(\mathcal S_0\)
\citep[Theorem~1.1, with characteristic zero and genus
zero]{FF_Guralnick2003}.
Every function field obtained by the operations in
Claim~\ref{clm:FFlatteskernel} is contained in a finite
compositum of Galois closures of genus-zero function-field
extensions. The Galois group of that compositum is a
subdirect product of the individual monodromy groups,
so it has no additional simple composition factors.
For completeness, this last assertion follows by induction
on the number of factors: the kernel of projection to
one factor projects normally into each remaining factor;
apply the induction hypothesis to that kernel and then
the composition series of the resulting extension.

Choose a nonalternating nonabelian finite simple group
\(S\notin\mathcal S_0\). Such groups exist, for
example among the infinitely many simple groups of
Lie type. The characteristic-zero fundamental-group
presentation and Riemann existence theorem give an
\(S\)-Galois cover of an open of
\(\mathbb P^1_{\overline{\mathbb Q}}\)
\citep[Expos\'es~XII, Th\'eor\`eme~5.1, and~XIII,
Corollaire~2.12]{FF_SGA1}.
It descends over a finite extension of the constants.
This cover cannot be dominated by an object of the
stated generated subcategory, since then \(S\)
would be a composition factor of one of the preceding
subdirect products. This is a scope restriction on the
closure argument. Noncofinality does not prove that the
remaining automorphism kernel is nontrivial.
